\subsection{Relatively connected maps}

DEFINITION 1. --- Let X and Y be topological spaces and let f be a continuous map from X to Y. The map f is said to be relatively connected if every point of Y has a fundamental system of neighborhoods consisting of sets V such that \(f^{-1}(V)\) is connected and nonempty.

Let \(f \colon X \to Y\) be a continuous map; for \(f\) to be relatively connected, it is necessary and sufficient that every point of \(Y\) has a neighborhood \(V\) such that the map \(f_V \colon f^{-1}(V) \to V\) deduced from \(f\) be relatively connected.

Let X and Y be topological spaces, and let \(f: X \to Y\) be a relatively connected continuous map. The image of f is dense in Y. For every open subset V of Y, the map \(f_{\mathrm{V}}: f^{-1}(\mathrm{V}) \to \mathrm{V}\) is relatively connected.

PROPOSITION 4. --- Let X and Y be topological spaces and let \(f: X \to Y\) be a relatively connected continuous map.

\begin{enumerate}
\def\labelenumi{\alph{enumi})}
\item
  For every connected component U of X, \(f(\overline{\mathrm{U}})\) is a connected, open, and closed subset of Y, and we have \(f^{-1}(\overline{f(\mathrm{U})}) = \mathrm{U}\) .
\item
  For every connected component V of Y, there exists a connected component U of X such that \(V = \overline{f(U)}\) . The map from U to V induced by f by passing to subsets is relatively connected.
\item
  The connected components of X (resp. of Y) are open and closed.
\item
Let V be the set of \(y \in Y\) which possess a neighborhood W such that \(f^{-1}(W)\) is connected and meets U; it is an open subset of Y. It contains \(\overline{f(U)}\) since f is relatively connected. Conversely, let \(y \in V\); let W be a neighborhood of \(y\) such that \(f^{-1}(W)\) is connected and meets U. For every neighborhood \(W'\) of \(y\) such that \(f^{-1}(W')\) is connected, \(f^{-1}(W \cap W')\) is not empty, therefore \(f^{-1}(W \cup W')\) is connected (TG, I, p. 81, prop. 2). By hypothesis, \(f^{-1}(W \cup W')\) meets U; one therefore has \(f^{-1}(W \cup W') \subset U\) and, a fortiori, \(W' \cap f(U) \neq \varnothing\). This proves that \(y \in \overline{f(U)}\); therefore, \(V = \overline{f(U)}\). In particular, the set \(\overline{f(U)}\) is open and closed in Y; moreover, prop. 1 (TG, I, p. 81) implies that it is connected.
\end{enumerate}

The preceding arguments also show that \(f^{-1}(\overline{f(\mathrm{U})}) \subset \mathrm{U}\) , whence the equality \(f^{-1}(\overline{f(\mathrm{U})}) = \mathrm{U}\) , the other inclusion being evident. In particular, U is open and closed in X.

\begin{enumerate}
\def\labelenumi{\alph{enumi})}
\setcounter{enumi}{1}
\tightlist
\item
Let V be the connected component of a point y of Y, let W be a neighborhood of y such that \(f^{-1}(W)\) is connected and nonempty, and let U be the connected component of X that contains \(f^{-1}(W)\). Since \(y \in \overline{f(U)}\), it follows from a) and from the definition of the connected component of y that \(\mathrm{V} = \overline{f(\mathrm{U})}\). Consequently, V is open and closed in Y.
\end{enumerate}

COROLLARY 1. --- Let X and Y be topological spaces and let \(f: X \to Y\) be a continuous and relatively connected map. By passing to connected components, the map \(f\) induces a bijection from the set of connected components of X onto the set of connected components of Y.

COROLLARY 2. --- Let X and Y be topological spaces and let \(f: X \to Y\) be a continuous map. For \(f\) to be relatively connected, it is necessary and sufficient that the following three properties be satisfied:

\begin{enumerate}
\def\labelenumi{\alph{enumi})}
\item
  The image f(X) is dense in Y;
\item
  The space Y is locally connected;
\item
  For every open and connected set V in Y, the set \(f(\mathrm{V})\) is connected.
\end{enumerate}

Suppose that the map f is relatively connected. The density of \(f(\mathrm{X})\) in Y follows from the definition of a relatively connected map. Let V be an open subset of Y; the map \(f_{\mathrm{V}}: f^{-1}(\mathrm{V}) \to \mathrm{V}\) is relatively connected. By prop. 4, the connected components of V are open and closed in V. It follows that Y is locally connected (TG, I, p. 85, prop. 11). To prove assertion c), it suffices to prove that X is connected if Y is. By the lemma, there exists a connected component U of X such that \(Y = \overline{f(U)}\) and \(f^{-1}(\overline{f(U)}) = U\) , whence U = X and X is connected.

Conversely, suppose that the conditions \(a\), \(b\), \(c\) are satisfied and let us show that f is relatively connected. Let \(y\) be a point of Y; since Y is locally connected, \(y\) admits a fundamental system of connected open neighborhoods. If W is such a neighborhood, the conditions \(c\) and \(a\) imply that \(f^{-1}(W)\) is connected and nonempty. Consequently, the map f is relatively connected.

COROLLARY 3. --- Let X and Y be topological spaces and let \(f\colon X\to Y\) be a continuous and relatively connected map. Let F be a set and let \(g\colon X\to F\) be a locally constant map. There exists a unique locally constant map \(h\colon Y\to F\) such that \(g=h\circ f\) .

The restriction of g to every connected component of X is constant. By Corollary 1, there exists a map \(h: Y \to F\) , constant on each connected component of Y, such that \(g = h \circ f\) . The map h is locally constant, because the connected components of Y are open. The uniqueness of such a map follows from the fact that \(f(X)\) is dense in Y.

Examples. --- 1) Let X be a topological space and let R be an equivalence relation on X. Denote by Y the quotient topological space X/R and by \(f\colon X \to Y\) the canonical map. Suppose that the equivalence classes of R are connected. Then, for every open and connected subset V of Y, the set \(f^{-1}(V)\) is connected (TG, I, p. 23, corollary 1 and p. 82, proposition 7). If the space Y is locally connected, the map f is thus relatively connected.

\begin{enumerate}
\def\labelenumi{\arabic{enumi})}
\setcounter{enumi}{1}
\tightlist
\item
  Let Y be a locally path-connected topological space and let X be an open subset of Y. The space \(\mathcal{C}_{c}(\mathbf{I};\mathbf{X})\) of paths in X is identified with a subspace of the space \(\mathcal{C}_{c}(\mathbf{I};\mathbf{Y})\) of paths in Y; suppose it is dense. The canonical injection of X into Y is then a relatively connected map.
\end{enumerate}

It suffices indeed to show that, for every nonempty connected open subset V of Y, the set \(V \cap X\) is connected and nonempty. The space \(\mathcal{C}_{c}(\mathbf{I};\mathrm{V})\) is a nonempty open set of \(\mathcal{C}_{c}(\mathbf{I};\mathrm{Y})\) ; it meets \(\mathcal{C}_{c}(\mathbf{I};\mathrm{X})\) , which proves that \(V \cap X \neq \varnothing\) . Let x and \(x'\) be points of \(V \cap X\) . Since \(V \cap X\) is locally path-connected, the points x and \(x'\) have open neighborhoods U and \(U'\) path-connected and contained in \(V \cap X\) . Since V is path-connected (III, p. 260, prop. 7), there exists a path in V joining x to \(x'\) . By the density hypothesis and the definition of the compact-convergence topology, there exists a path in \(V \cap X\) joining a point of U to a point of \(U'\) . There then exists a path joining x to \(x'\) in \(V \cap X\) , which proves that the set \(V \cap X\) is connected.

PROPOSITION 5. --- Let X and Y be topological spaces and let \(f: X \to Y\) be a continuous and relatively connected map. For every pair \((T, T')\) of coverings of Y, the map \(f^{*}: \mathcal{C}_{Y}(T; T') \to \mathcal{C}_{X}(X \times_{Y} T; X \times_{Y} T')\) is bijective.

Let \(\mathcal{F}\) be the sheaf on X of X-morphisms from \(X \times_{Y} T\) into \(X \times_{Y} T'\) and let \(\mathcal{G}\) be the sheaf on Y of Y-morphisms from T into \(T'\) (I, p. 45, example 4). For every open set U of Y, set \(\varphi_{\mathrm{U}} = (f_{\mathrm{U}})^{*}: \mathcal{G}(\mathrm{U}) \to \mathcal{F}(f^{-1}(\mathrm{U}))\) . The maps \(\varphi_{\mathrm{U}}\) define a morphism of sheaves \(\varphi: \mathcal{G} \to \varphi_{*}(\mathcal{F})\) and it suffices to prove that \(\varphi\) is an isomorphism of sheaves.

Since Y is locally connected (IV, p. 354, cor. 2), the connected open sets over which T and T' are trivializable form a base of the topology of Y. By corollary 2 of I, p. 55, it suffices to prove that for such an open set U, the map \(\varphi_{\mathrm{U}}\) is bijective, which allows us to assume that Y is connected and that the coverings T and T' are the trivial coverings \(Y \times F\) and \(Y \times F'\) where F and F' are sets equipped with the discrete topology. The map \((x, (y, t)) \mapsto (x, t)\) identifies the X-space \(X \times_{Y} (Y \times F)\) with \(X \times F\) (resp. the X-space \(X \times_{Y} (Y \times F')\) with \(X \times F'\)). Since the space X is connected (IV, p. 354, cor. 2), the sets \(\mathcal{C}_{\mathrm{Y}}(Y \times F; Y \times F')\) and \(\mathcal{C}_{\mathrm{X}}(X \times F; X \times F')\) are both identified with the set \(\mathcal{F}(F;F')\) of maps from F into F', and the map \(f^{*}\) is identified with the identity map of \(\mathcal{F}(F;F')\). This concludes the proof.

COROLLARY 1. --- Let X and Y be topological spaces and let \(f: X \to Y\) be a relatively connected continuous map. Let T and T' be coverings of Y. If the coverings \(X \times_Y T\) and \(X \times_Y T'\) of X are isomorphic, then the coverings T and T' are isomorphic.

Indeed, let \(h \colon X \times_{Y} T \to X \times_{Y} T'\) and \(h' \colon X \times_{Y} T' \to X \times_{Y} T\) be X-isomorphisms reciprocal to each other. By Proposition 5, there exist Y-morphisms \(g \colon \mathrm{T} \to \mathrm{T}'\) and \(g' \colon \mathrm{T}' \to \mathrm{T}\) such that \(f^{*}(g) = h\) and \(f^{*}(g') = h'\) . One then has \(f^{*}(g' \circ g) = f^{*}(g') \circ f^{*}(g) = \mathrm{Id}_{\mathrm{X} \times_{\mathrm{Y}} \mathrm{T}}\) , therefore \(g' \circ g = \mathrm{Id}_{\mathrm{T}}\) , because the map \(f^{*}\) is injective. Similarly, \(g \circ g' = \mathrm{Id}_{\mathrm{T}'}\) . The coverings T and \(\mathrm{T}'\) are therefore isomorphic.

COROLLARY 2. --- Let X and Y be topological spaces and let \(f: X \to Y\) be a continuous and relatively connected map. If the space X is simply connected, then so is the space Y.

By Corollary 1, every covering of Y is indeed trivializable.

COROLLARY 3. --- Let X be a locally path-connected topological space, Y a loop-disentanglable topological space, and let \(f: X \to Y\) be a relatively connected continuous map. For every point \(x\) of X, the homomorphism \(\pi_1(f, x): \pi_1(X, x) \to \pi_1(Y, f(x))\) is surjective.

By IV, p. 353, prop. 4, we may assume that the spaces X and Y are connected. The corollary then follows from proposition 5 and proposition 2 of IV, p. 352.

PROPOSITION 6. --- Let X and Y be topological spaces, and let \(f: \mathrm{X} \to \mathrm{Y}\) be a continuous and relatively connected map. Suppose that every point of Y has an open neighborhood V whose inverse image \(f^{-1}(V)\) is simply connected. Then, for every covering Z of X, there exists a covering T of Y such that \(X \times_{Y} T\) is X-isomorphic to Z.

Let \(\mathcal{U}\) be the set of open subsets of Y whose inverse image in X is simply connected. For every \(V \in \mathcal{U}\), the covering \(f^{-1}(V) \times_{X} Z\) of \(f^{-1}(V)\) is trivializable; there thus exists a discrete space \(F_V\) and an isomorphism of coverings \(g_V: f^{-1}(V) \times_{X} Z \to f^{-1}(V) \times F_V\) . For every pair \((V, V')\) of open sets belonging to \(\mathcal{U}\), the map \(f_{V \cap V'}: f^{-1}(V \cap V') \to V \cap V'\) is relatively connected. By Proposition 5 of IV, p. 356, there exists a unique isomorphism of coverings of \(V \cap V'\), \(h_{V',V}: (V \cap V') \times F_V \to (V \cap V') \times F_{V'}\) such that one has \(f^*(h_{V',V})(x,t) = g_{V'}(g_V^{-1}(x,t))\) for all \(x \in V \cap V'\) and every \(t \in F_V\). If V, \(V'\), \(V''\) are elements of \(\mathcal{U}\), we have \(h_{V'',V}(x,t) = h_{V'',V'}(h_{V',V}(x,t))\) for all \(x \in V \cap V' \cap V''\) and every \(t \in F_V\). There then exists a unique Y-space T and, for every \(V \in \mathcal{U}\), an isomorphism \(h_V: T_V \to V \times F_V\) such that one has \(h_{V',V}(x,t) = h_{V'} \circ h_V^{-1}(x,t)\) for every pair \((V, V')\) of open sets belonging to \(\mathcal{U}\), all \(x \in V \cap V'\) and every \(t \in F_V\) (cf. TG, I, p. 16).

The space T is in particular a covering of Y. Moreover, there exists a unique map from \(X \times_{Y} T\) over Z whose restriction to \(f^{-1}(V) \times_{Y} T\) is given by \(g_{V}^{-1} \circ f^{*}(h_{V})\) and it is an isomorphism of X-spaces, whence the proposition.

COROLLARY. --- Let X and Y be loop-disentanglable topological spaces and let \(f: X \to Y\) be a continuous and relatively connected map. Suppose that every point of Y has a neighborhood V whose inverse image \(f^{-1}(V)\) is simply connected. Then, for every point x in X, the homomorphism \(\pi_{1}(f,x): \pi_{1}(X,x) \to \pi_{1}(Y,f(x))\) is bijective.

One may assume the spaces X and Y are connected (IV, p. 353, prop. 4). By Corollary 3 (IV, p. 357), the homomorphism \(\pi_1(f,x)\) is surjective. Proposition 6 and Proposition 1 of IV, p. 351 imply that it is injective.

Remark. --- Let Y be a topological space, and let X, X', and Y' be subspaces of Y such that X ⊂ X' ⊂ Y' ⊂ Y. Suppose that the canonical injection of X into Y is relatively connected. For every connected open subset V of Y, the set V ∩ X is connected and dense in V (IV, p. 354, cor. 2); consequently, the set V ∩ X' is connected (TG, I, p. 81, prop. 1). This proves that the canonical injection of X' into Y' is relatively connected.

Let V be an open subset of Y; by the preceding, the canonical injection of \(V \cap X\) into \(V \cap X'\) is relatively connected. If the set \(V \cap X\) is simply connected, \(V \cap X'\) is also, by virtue of Corollary 2 of IV, p. 357. Consequently, if the canonical injection of X into Y satisfies the hypotheses of Proposition 6, then so does the canonical injection of \(X'\) into \(Y'\) .

Example. --- Let Y be a locally finite-dimensional differentiable manifold and let Z be a closed submanifold of Y (VAR, R, 5.8.3). Set X = Y - Z.

\begin{enumerate}
\def\labelenumi{\alph{enumi})}
\item
  If the codimension of Z is at least 2 at every point, the canonical injection of X into Y is relatively connected. Indeed, let z be a point of Z; there exists an open neighborhood V of z in Y and a homeomorphism \(\varphi\) of V onto a finite-dimensional vector space E over R such that \(\varphi(V \cap Z)\) is a vector subspace F of E whose codimension is \(\geqslant 2\) . The set E-F is connected (TG, VI, p. 5, proposition 4), whence the assertion.
\item
  Suppose moreover that the codimension of Z is at least 3 at every point; the hypotheses of Proposition 6 and its corollary are then satisfied, since, with the notation of the preceding paragraph, E-F is simply connected (I, p. 128, example 4).
\end{enumerate}

Since differentiable manifolds are loop-disentanglable spaces (IV, p. 347), the results of this section admit the following particular case.

Let Y be a locally finite-dimensional differentiable manifold, Z a closed submanifold of Y, and i the canonical injection of Y−Z into Y.

\begin{enumerate}
\def\labelenumi{\alph{enumi})}
\item
  If the codimension of Z in Y is ≥ 1 at every point, then the variety Y−Z is dense in Y and the map \(\pi_{0}(i)\) is surjective.
\item
  If the codimension of Z in Y is ≥ 2 at every point, the map \(\pi_{0}(i)\) is bijective and, for every point x in Y-Z, the map \(\pi_{1}(i,x)\) is surjective.
\item
  If the codimension of Z in Y is \(\geqslant 3\) at every point, the maps \(\pi_{0}(i)\) and \(\pi_{1}(i,x)\) are bijective, for every point x in Y-Z.
\end{enumerate}

